\clearpage\thispagestyle{empty}
\pdfbookmark[0]{Sources for the review of Harding’s star atlas}{harding1809-sources}
\begin{center}\textbf{Sources for the review of Harding’s star atlas}\end{center}
\noindent This AI-translated review follows Gauss’s text on pp.~510--512 of \href{https://gdz.sub.uni-goettingen.de/id/PPN236007467}{\textit{Werke}, vol.~VI} (Göttingen, 1874). It appeared unsigned under the heading \emph{Göttingen} in \textit{Göttingische gelehrte Anzeigen}, no.~61, 17 April 1809, pp.~\href{https://images.p.sub.uni-goettingen.de/iiif/image/gdz:PPN319721507_1809:00000603/full/1750,/0/default.jpg}{601}--\href{https://images.p.sub.uni-goettingen.de/iiif/image/gdz:PPN319721507_1809:00000609/full/1750,/0/default.jpg}{607}. The digitized copy is held by the Göttingen State and University Library.
\par\smallskip\noindent The first printing calls the final chart in the list IX. The collected edition reads IV there, while listing I, II, V and IX at the start. Both printings give the right ascension endpoint 81 without a degree sign. The translation follows the collected edition.
\par\smallskip\noindent Gauss refers to his August 1804 paper \hyperlink{grenzen1804.1}{\emph{On the Limits of the Geocentric Positions of the Planets}}. Delambre’s quotation from Plutarch means \emph{a deed owing to merit favoured by fortune, rather than to chance}. The word order differs from \href{https://penelope.uchicago.edu/Thayer/E/Roman/Texts/Plutarch/Lives/Timoleon*.html#36.4}{\textit{Life of Timoleon}, 36.4}. \href{https://www.davidrumsey.com/luna/servlet/detail/RUMSEY~8~1~372247~90138986:IX-Atlas-novus-coelestis-XXVII-tabu}{Harding’s chart IX} in the David Rumsey Historical Map Collection belongs to the 1822 atlas of 27 sheets. Translation and typesetting produced by AI.
\clearpage
